\chapter{Trigonométrie dans le triangle rectangle}
\label{exercices:15}
\section*{Exercice 15.1 -- Les trois rapports}
$ABC$ est rectangle en $A$, avec $AB=8$, $AC=6$, $BC=10$. Pour $\alpha=\widehat{ABC}$, calculer sinus, cosinus et tangente, puis $\alpha$ au degré près.

\section*{Exercice 15.2 -- Hauteur et recul}
Une échelle de $4\unite m$ forme un angle de $70\degre$ avec un sol horizontal et touche un mur vertical. Calculer la hauteur atteinte et la distance du pied au mur, au centimètre près.

\section*{Exercice 15.3 -- Visée d'un sommet}
À $15\unite m$ horizontalement d'un bâtiment, un observateur dont les yeux sont à $1{,}50\unite m$ du sol vise le sommet sous un angle de $28\degre$. Calculer la hauteur du bâtiment. Expliciter les hypothèses géométriques.

\section*{Exercice 15.4 -- Identité et pente}
Démontrer $\sin^2\alpha+\cos^2\alpha=1$ dans un triangle rectangle. Une pente de $24\,\%$ fait-elle un angle de $24\degre$ ? Calculer son angle au dixième près.
